\documentclass[a4paper,12pt]{article} % Dokumentenklasse
\usepackage[ngerman]{babel} % Deutsche Spracheinstellungen
\usepackage[utf8]{inputenc} % UTF-8 Eingabekodierung
\usepackage[T1]{fontenc} % T1 Schriftkodierung
\usepackage[a4paper, left=2cm, right=2cm, top=2cm, bottom=2cm]{geometry} % Seitenränder
\usepackage{amsmath} % Für erweiterte mathematische Umgebungen
\usepackage{amsfonts} % Für erweiterte mathematische Schriftarten
\usepackage{amssymb} % Für erweiterte mathematische Symbole
\usepackage{graphicx} % Für Grafiken
\usepackage{array} % Für bessere Tabellenanpassungen
\usepackage{caption} % Für Tabellenunterschriften
\usepackage{array} % Für bessere Tabellenanpassungen
\usepackage{enumitem} % Für benutzerdefinierte Aufzählungen

\renewcommand{\thesubsubsection}{\Roman{subsubsection}.} % Unterunterabschnittsnummerierung

\begin{document}


\title{Elementare Gleichungen - Aufgaben} % Titel des Dokuments
\author{Alexander Roux - bearbeitet von Jakob Janzen} % Autor des Dokuments
\date{\today} % Datum des Dokuments
\maketitle % Titelseite erstellen

\pagenumbering{Roman} % Römische Seitenzahlen für das Inhaltsverzeichnis
\tableofcontents % Inhaltsverzeichnis erstellen

\newpage % Neue Seite nach dem Inhaltsverzeichnis
\pagenumbering{arabic} % Arabische Seitenzahlen für den Hauptteil

\section{Einführung}

\subsection{Aufgaben}

\subsubsection{Löse die folgenden Glleichungen:} % I.

\begin{enumerate}

  % 1.
  \item Löse die Gleichung \(2x + 4 = 18\).
  \[
  \begin{aligned}
    2x + 4 &= 18 \quad &\big| -4 \\
    2x &= 14 \quad &\big| \div 2 \\
    x &= \mathbf{7}
  \end{aligned}
  \]

  % 2.
  \item Löse die Gleichung \(6x + 3 = 33\).
  \[
  \begin{aligned}
    6x + 3 &= 33 &\quad\big| -3 \\
    6x &= 30 &\quad\big| \div 6 \\
    x &= \mathbf{5}
  \end{aligned}
  \]

  % 3.
  \item Löse die Gleichung \(5x + 8 = 13\).
  \[
  \begin{aligned}
    5x + 8 &= 13 &\quad\big| -8 \\
    5x &= 5 &\quad\big| \div 5 \\
    x &= \mathbf{1}
  \end{aligned}
  \]

  % 4.
  \item Löse die Gleichung \(3x + 9 = 18\).
  \[
  \begin{aligned}
    3x + 9 &= 18 &\quad\big| -9 \\
    3x &= 9 &\quad\big| \div 3 \\
    x &= \mathbf{3}
  \end{aligned}
  \]

  % 5.
  \item Löse die Gleichung \(6x + 9 = 45\).
  \[
  \begin{aligned}
    6x + 9 &= 45 &\quad\big| -9 \\
    6x &= 36 &\quad\big| \div 6 \\
    x &= \mathbf{6}
  \end{aligned}
  \]

  % 6.
  \item Löse die Gleichung \(4x + 9 = 29\).
  \[
  \begin{aligned}
    4x + 9 &= 29 &\quad\big| -9 \\
    4x &= 20 &\quad\big| \div 4 \\
    x &= \mathbf{5}
  \end{aligned}
  \]

  % 7.
  \item Löse die Gleichung \(7x + 4 = 25\).
  \[
  \begin{aligned}
    7x + 4 &= 25 &\quad\big| -4 \\
    7x &= 21 &\quad\big| \div 7 \\
    x &= \mathbf{3}
  \end{aligned}
  \]

  % 8.
  \item Löse die Gleichung \(2x + 7 = 11\).
  \[
  \begin{aligned}
    2x + 7 &= 11 &\quad\big| -7 \\
    2x &= 4 &\quad\big| \div 2 \\
    x &= \mathbf{2}
  \end{aligned}
  \]

  % 9.
  \item Löse die Gleichung \(4x + 8 = 36\).
  \[
  \begin{aligned}
    4x + 8 &= 36 &\quad\big| -8 \\
    4x &= 28 &\quad\big| \div 4 \\
    x &= \mathbf{7}
  \end{aligned}
  \]

  % 10.
  \item Löse die Gleichung \(6x + 4 = 46\).
  \[
  \begin{aligned}
    6x + 4 &= 46 &\quad\big| -4 \\
    6x &= 42 &\quad\big| \div 6 \\
    x &= \mathbf{7}
  \end{aligned}
  \]

\end{enumerate}

\section{Negative Zahlen}

\subsection{Aufgaben zu Kapitel 2.1 bis 2.3 \\
  - Was sind negative Zahlen? \\
  - Negative Zahlen im Alltag \\
  - Addition mit negativen Zahlen}

\subsubsection{Berechne die folgenden Summen:} % I.

\begin{enumerate}

  \item \(4 + (-3) = 4 - 3 = \mathbf{1}\) % 1.

  \item \(5 + (-2) = 5 - 2 = \mathbf{3}\) % 2.

  \item \(4 + (-6) = 4 - 6 = \mathbf{-2}\) % 3.

  \item \(9 + (-7) = 9 - 7 = \mathbf{5}\) % 4.

  \item \(8 + (-3) = 8 - 3 = \mathbf{5}\) % 5.

  \item \(4 + (-17) = 4 - 17 = \mathbf{-13}\) % 6.

  \item \(14 + (-5) = 14 - 5 = \mathbf{9}\) % 7.

  \item \(1 + (-8) = 1 - 8 = \mathbf{-7}\) % 8.

  \item \(9 + (-12) = 9 - 12 = \mathbf{-3}\) % 9.

  \item \(7 + (-20) = 7 - 20 = \mathbf{-13}\) % 10.

  \item \(23 + (-12) = 23 - 12 = \mathbf{11}\) % 11.

  \item \(33 + (-21) = 33 - 21 = \mathbf{12}\) % 12.

  \item \(14 + (-6) = 14 - 6 = \mathbf{8}\) % 13.

  \item \(25 + (-30) = 25 - 30 = \mathbf{-5}\) % 14.

  \item \(5 + (-9) = 5 - 9 = \mathbf{-4}\) % 15.

  \item \(17 + (-8) = 17 - 8 = \mathbf{9}\) % 16.

  \item \(11 + (-18) = 11 - 18 = \mathbf{-7}\) % 17.

  \item \(17 + (-25) = 17 - 25 = \mathbf{-8}\) % 18.

  \item \(27 + (-6) = 27 - 6 = \mathbf{21}\) % 19.

  \item \(13 + (-43) = 13 - 43 = \mathbf{-30}\) % 20.

  \item \(4 + (-1) = 4 - 1 = \mathbf{3}\) % 21.

  \item \(-2 + 3 = \mathbf{1}\) % 22.

  \item \(-2 + (-4) = -2 - 4 = \mathbf{-6}\) % 23.

  \item \(3 + (-5) = 3 - 5 = \mathbf{-2}\) % 24.

  \item \(-2 + 1 = \mathbf{-1}\) % 25.

  \item \(-5 + (-3) = -5 - 3 = \mathbf{-8}\) % 26.

  \item \(4 + (-5) = 4 - 5 = \mathbf{-1}\) % 27.

  \item \(-3 + (-4) = -3 - 4 = \mathbf{-7}\) % 28.

  \item \(-2 + (-7) = -2 - 7 = \mathbf{-9}\) % 29.

  \item \(2 + (-8) = 2 - 8 = \mathbf{-6}\) % 30.

  \item \(-1 + (-1) = -1 - 1 = \mathbf{-2}\) % 31.

  \item \(4 + (-7) = 4 - 7 = \mathbf{-3}\) % 32.

  \item \(-2 + (-3) = -2 - 3 = \mathbf{-5}\) % 33.

  \item \(-3 + 5 = \mathbf{2}\) % 34.

  \item \(-3 + (-4) = -3 - 4 = \mathbf{-7}\) % 35.

  \item \(-1 + (-7) = -1 - 7 = \mathbf{-8}\) % 36.

  \item \(-3 + 4 = \mathbf{1}\) % 37.

  \item \(2 + (-6) = 2 - 6 = \mathbf{-4}\) % 38.

  \item \(-3 + 2 = \mathbf{-1}\) % 39.

  \item \(-2 + (-8) = -2 - 8 = \mathbf{-10}\) % 40.

  \item \(2 + (-2) = 2 - 2 = \mathbf{0}\) % 41.

  \item \(-1 + 1 = \mathbf{0}\) % 42.

  \item \(-3 + 3 = \mathbf{0}\) % 43.

  \item \(5 + (-5) = 5 - 5 = \mathbf{0}\) % 44.

  \item \(-7 + 7 = \mathbf{0}\) % 45.

  \item \(13 + (-13) = 13 - 13 = \mathbf{0}\) % 46.

  \item \(-8 + 8 = \mathbf{0}\) % 47.

  \item \(35 + (-35) = 35 - 35 = \mathbf{0}\) % 48.

\end{enumerate}

\subsubsection{Addiere die benachbarten Zahlen und schreibe das Ergebnis darüber.
Mithilfe der obersten Zahl kann man erkennen, ob man richtig gerechnet hat.} % II.

\begin{enumerate}[label=\alph*)]

  \item % a)
  \vspace{1cm}
  \begin{center}
    \begin{tabular}{c c c c c c c c c}
      & & & & -6 & & & & \\
      & & & -4 & & -2 & & & \\
      & & -1 & & -3 & & 1 & & \\
      & 1 & & -2 & & -1 & & 2 & \\
      -2 & & 3 & & -5 & & 4 & & -2 \\
    \end{tabular}
  \end{center}

  \item % b)
  \vspace{1cm}
  \begin{center}
    \begin{tabular}{c c c c c c c c c}
      & & & & -8 & & & & \\
      & & & -1 & & -7 & & & \\
      & & 1 & & -2 & & -5 & \\
      & 1,5 & & -0,5 & & -1,5 & & -3,5 & \\
      3,5 & & -2 & & 1,5 & & -3 & & -0,5 \\
    \end{tabular}
  \end{center}

\end{enumerate}
\vspace{1cm}

\subsection{Aufgaben zu Kapitel 2.4 \\
  - Substraktion mit negativen Zahlen}

\subsubsection{Berechne die folgenden Differenzen:}

\begin{enumerate} % I.

  \item \(4 - 3 = \mathbf{1}\) % 1.

  \item \(8 - 5 = \mathbf{3}\) % 2.

  \item \(5 - 7 = \mathbf{-2}\) % 3.

  \item \(10 - 7 = \mathbf{3}\) % 4.

  \item \(6 - 19 = \mathbf{-13}\) % 5.

  \item \(13 - 5 = \mathbf{8}\) % 6.

  \item \(2 - 9 = \mathbf{-7}\) % 7.

  \item \(11 - 14 = \mathbf{-3}\) % 8.

  \item \(6 - 30 = \mathbf{-24}\) % 9.

  \item \(34 - 23 = \mathbf{11}\) % 10.

  \item \(15 - 7 = \mathbf{8}\) % 11.

  \item \(36 - 41 = \mathbf{-5}\) % 12.

  \item \(4 - 8 = \mathbf{-4}\) % 13.

  \item \(12 - 17 = \mathbf{-5}\) % 14.

  \item \(18 - 26 = \mathbf{-8}\) % 15.

  \item \(26 - 5 = \mathbf{21}\) % 16.

  \item \(15 - 45 = \mathbf{-30}\) % 17.

  \item \(-9 - (-2) = -9 + 2 = \mathbf{-7}\) % 18.

  \item \(3 - (-2) = 3 + 2 = \mathbf{5}\) % 19.

  \item \(-4 - (-1) = -4 + 1 = \mathbf{-3}\) % 20.

  \item \(-2 - 3 = \mathbf{-5}\) % 21.

  \item \(-2 - 4 = \mathbf{-6}\) % 22.

  \item \(3 - (-5) = 3 + 5 = \mathbf{8}\) % 23.

  \item \(-2 - 1 = \mathbf{-3}\) % 24.

  \item \(-5 - 3 = \mathbf{-8}\) % 25.

  \item \(-3 - (-4) = -3 + 4 = \mathbf{1}\) % 26.

  \item \(-2 - (-4) = -2 + 4 = \mathbf{2}\) % 27.

  \item \(-1 - (-1) = -1 + 1 = \mathbf{0}\) % 28.

  \item \(4 - (-7) = 4 + 7 = \mathbf{11}\) % 29.

  \item \(-2 - (-3) = -2 + 3 = \mathbf{1}\) % 30.

  \item \(-3 - 5 = \mathbf{-8}\) % 31.

  \item \(-3 - (-4) = -3 + 4 = \mathbf{1}\) % 32.

  \item \(-1 - 7 = \mathbf{-8}\) % 33.

  \item \(-3 - 4 = \mathbf{-7}\) % 34.

  \item \(2 - (-6) = 2 + 6 = \mathbf{8}\) % 35.

  \item \(-3 - 2 = \mathbf{-5}\) % 36.

  \item \(-2 - (-8) = -2 + 8 = \mathbf{6}\) % 37.

  \item \(-2 - 2 = \mathbf{-4}\) % 38.

  \item \(-1 - 1 = \mathbf{-2}\) % 39.

  \item \(4 - (-4) = 4 + 4 = \mathbf{8}\) % 40.

  \item \(7 - (-7) = 7 + 7 = \mathbf{14}\) % 41.

  \item \(35 - (-35) = 35 + 35 = \mathbf{70}\) % 42.

  \item \(3 - 5 = \mathbf{-2}\) % 43.

  \item \(-4 - 6 = \mathbf{-10}\) % 44.

  \item \(-9 - (-6) = -9 + 6 = \mathbf{-3}\) % 45.

  \item \(-4 - (-8) = -4 + 8 = \mathbf{4}\) % 46.

  \item \(5 - (-4) = 5 + 4 = \mathbf{9}\) % 47.

  \item \(4 - 7 = \mathbf{-3}\) % 48.

\end{enumerate}

\subsubsection{Subtrahiere die benachbarten Zahlen (linke Zahl minus rechte Zahl)
und schreibe das Ergebnis darüber. Mithilfe der obersten Zahl kann man erkennen,
ob man richtig gerechnet hat.} % II.

\begin{enumerate}[label=\alph*)]

  \item % a)
  \vspace{1cm}
  \begin{center}
    \begin{tabular}{c c c c c c c c c}
      & & & & -62 & & & & \\
      & & & -30 & & 32 & & & \\
      & & -13 & & 17 & & -15 & & \\
      & -5 & & 8 & & -9 & & 6 & \\
      -2 & & 3 & & -5 & & 4 & & -2 \\
    \end{tabular}
  \end{center}

  \item % b)
  \vspace{1cm}
  \begin{center}
    \begin{tabular}{c c c c c c c c c}
      & & & & 32 & & & & \\
      & & & 17 & & -15 & & & \\
      & & 9 & & -8 & & 7 & \\
      & 5,5 & & -3,5 & & 4,5 & & -2,5 & \\
      3,5 & & -2 & & 1,5 & & -3 & & -0,5 \\
    \end{tabular}
  \end{center}

\end{enumerate}
\vspace{1cm}

\subsection{Aufgaben zu Kapitel 2.5 \\
  - Verschmelzung von Addition und Subtraktion}

\subsubsection{Berechne die folgenden Summen/Differenzen:} % I.

\begin{enumerate}

  \item \(-27 + 8 = \mathbf{-19}\) % 1.

  \item \(-18 - 13 = \mathbf{-31}\) % 2.

  \item \(-28 - 1 = \mathbf{-29}\) % 3.

  \item \(-24 + 17 = \mathbf{-7}\) % 4.

  \item \(-7 - 12 = \mathbf{-19}\) % 5.

  \item \(-11 + 17 = \mathbf{6}\) % 6.

  \item \(-21 + 15 = \mathbf{-6}\) % 7.

  \item \(17 - 15 = \mathbf{2}\) % 8.

  \item \(-16 - 17 = \mathbf{-33}\) % 9.

  \item \(-29 + 7 = \mathbf{-22}\) % 10.

  \item \(-11 - 23 = \mathbf{-34}\) % 11.

  \item \(6 - 30 = \mathbf{-24}\) % 12.

  \item \(-5 + 2 = \mathbf{-3}\) % 13.

  \item \(2 - 11 = \mathbf{-9}\) % 14.

  \item \(11 - 26 = \mathbf{-15}\) % 15.

  \item \(-16 - 16 = \mathbf{-32}\) % 16.

  \item \(6 - 23 = \mathbf{-17}\) % 17.

  \item \(21 - 22 = \mathbf{-1}\) % 18.

  \item \(-28 - 13 = \mathbf{-41}\) % 19.

  \item \(29 - 22 = \mathbf{7}\) % 20.

  \item \(5 - 16 = \mathbf{-11}\) % 21.

  \item \(-24 - 11 = \mathbf{-35}\) % 22.

  \item \(-7 - 5 = \mathbf{-12}\) % 23.

  \item \(-13 - 26 = \mathbf{-39}\) % 24.

  \item \(-1 - 12 = \mathbf{-13}\) % 25.

  \item \(-3 - 5 = \mathbf{-8}\) % 26.

  \item \(-21 + 8 = \mathbf{-13}\) % 27.

  \item \(-16 + 12 = \mathbf{-4}\) % 28.

  \item \(-6 - 1 = \mathbf{-7}\) % 29.

  \item \(-3 + 7 = \mathbf{4}\) % 30.

\end{enumerate}

\subsection{Aufgaben zu Kapitel 2.6 bis 2.8 \\
  - Addition und Subtraktion mehrerer Zahlen \\
  - Klammern subtrahieren und addieren \\
  - Beispiele zur Addition und Subtraktion}

\subsubsection{Berechne die folgenden Summen/Differenzen:} % I.

\begin{enumerate}

  \item \(-1 + 5 - 3 = 4 - 3 = \mathbf{1}\) % 1.

  \item \(1 - 5 - 3 = -4 - 3 = \mathbf{-7}\) % 2.

  \item \(1 + 5 - 3 = 6 - 3 = \mathbf{3}\) % 3.

  \item \(-1 - 5 - 3 = -6 - 3 \mathbf{-9}\) % 4.

  \item \(-1 + 5 + 3 = 4 + 3 = \mathbf{7}\) % 5.

  \item \(1 - 5 + 3 = -4 + 3 = \mathbf{-1}\) % 6.

  \item \(1 + 5 + 3 = 6 + 3 = \mathbf{9}\) % 7.

  \item \(-1 - 5 + 3 = -6 + 3 = \mathbf{-3}\) % 8.

  \item \(-2 + 7 + 4 = 5 + 4 = \mathbf{9}\) % 9.

  \item \(-4 + 2 - 7 = -2 - 7 = \mathbf{-9}\) % 10.

  \item \(-7 + 4 - 2 = -3 - 2 = \mathbf{-5}\) % 11.

  \item \(-5 - 6 - 2 = -11 - 2 = \mathbf{-13}\) % 12.

  \item \(2 - 6 - 3 = -4 - 3 = \mathbf{-7}\) % 13.

  \item \(-4 - 1 + 4 = -5 + 4 = \mathbf{-1}\) % 14.

  \item \(7 + 3 - 2 = 10 - 2 = \mathbf{8}\) % 15.

  \item \(-2 + 7 + 1 = 5 + 1 = \mathbf{6}\) % 16.

  \item \(1 - 4 + 5 = -3 + 5 = \mathbf{2}\) % 17.

  \item \(-7 + 2 - 9 = -5 - 9 = \mathbf{-14}\) % 18.

  \item \(-2 - 5 - 6 = -7 - 6 = \mathbf{-13}\) % 19.

  \item \(1 - 5 - 8 = -4 - 8 = \mathbf{-12}\) % 20.

  \item \(8 - 5 + 7 - 20 = 3 + 7 - 20 = 10 - 20 = \mathbf{-10}\) % 21.

  \item \(-9 - 4 + 11 - 17 = -13 + 11 - 17 = -2 - 17 = \mathbf{-19}\) % 22.

\end{enumerate}

\subsubsection{Berechne die folgenden Summen/Differenzen:} % II.

\begin{enumerate}

  \item \(-16 - 6 + 19 = -22 + 19 = \mathbf{-3}\) % 1.

  \item \(-5 + 15 + 2 - 13 + 9 = 10 + 2 - 13 + 9 = 12 - 13 + 9 = -1 + 9 = \mathbf{8}\) % 2.

  \item \(-4 - 10 - 18 - 13 + 8 = -14 - 18 - 13 + 8 = -32 - 13 + 8 = -45 + 8 = \mathbf{-37}\) % 3.

  \item \(-18 - 6 - 7 - 9 + 15 = -24 - 7 - 9 + 15 = -31 - 9 + 15 = -40 + 15 = \mathbf{-25}\) % 4.

  \item \(-17 + 14 - 9 = -3 - 9 = \mathbf{-12}\) % 5.

  \item \(-20 - 1 + 3 = -21 + 3 = \mathbf{-18}\) % 6.

  \item \(16 - 12 + 13 - 29 = 4 + 13 - 29 = 17 - 29 = \mathbf{-12}\) % 7.

  \item \(-2 - 29 + 6 = -31 + 6 = \mathbf{-25}\) % 8.

  \item \(-17 + 14 - 9 = -3 - 9 = \mathbf{-12}\) % 9.

  \item \(-10 - 3 - 27 = -13 - 27 = \mathbf{-40}\) % 10.

  \item \(-19 - 13 + 5 - 1 - 24 = -32 + 5 - 1 - 24 = -27 - 1 - 24 = -28 - 24 = \mathbf{-52}\) % 11.

  \item \(-21 - 2 + 11 = -23 + 11 = \mathbf{-12}\) % 12.

  \item \(-25 - 4 + 13 = -29 + 13 = \mathbf{-16}\) % 13.

  \item \(-1 - 17 + 20 - 14 = -18 + 20 - 14 = 2 - 14 = \mathbf{-12}\) % 14.

  \item \(12 + 27 - 3 = 39 - 3 = \mathbf{36}\) % 15.

  \item \(23 - 26 - 14 + 8 + 27 = -3 - 14 + 8 + 27 = -17 + 8 + 27 = -9 + 27 = \mathbf{18}\) % 16.

  \item \(26 - 24 + 16 - 5 - 10 = 2 + 16 - 5 - 10 = 18 - 5 - 10 = 13 - 10 = \mathbf{3}\) % 17.

  \item \(-6 + 13 - 11 = 7 - 11 = \mathbf{-4}\) % 18.

  \item \(17 - 7 + 8 = 10 + 8 = \mathbf{18}\) % 19.

  \item \(-21 + 3 + 5 + 16 + 10 = -18 + 5 + 16 + 10 = -13 + 16 + 10 = 3 + 10 = \mathbf{13}\) % 20.

  \item \(14 - 1 + 10 - 22 - 1 = 13 + 10 - 22 - 1 = 23 - 22 - 1 = 1 - 1 = \mathbf{0}\) % 21.

  \item \(-6 - 11 + 17 - 12 + 13 = -17 + 17 - 12 + 13 = -12 + 13 = \mathbf{1}\) % 22.

  \item \(6 + 29 - 13 = 35 - 13 = \mathbf{22}\) % 23.

  \item \(-16 - 23 - 25 = -39 - 25 = \mathbf{-64}\) % 24.

  \item \(-18 + 9 - 20 = -9 - 20 = \mathbf{-29}\) % 25.

  \item \(-13 + 22 + 7 + 21 = 9 + 7 + 21 = 16 + 21 = \mathbf{37}\) % 26.

  \item \(1 + 10 - 22 + 6 - 11 = 11 - 22 + 6 - 11 = -11 + 6 - 11 = -5 - 11 = \mathbf{-16}\) % 27.

  \item \(16 - 17 + 25 = -1 + 25 = \mathbf{24}\) % 28.

  \item \(-5 - 29 - 8 - 23 = -34 - 8 - 23 = -42 - 23 = \mathbf{-65}\) % 29.

  \item \(-13 + 4 - 12 - 7 = -9 - 12 - 7 = -21 - 7 = \mathbf{-28}\) % 30.

\end{enumerate}

\subsubsection{Berechne:} % III.

\begin{enumerate}

  \item \(8 - (9 - 4) = 8 - 5 = \mathbf{3}\) % 1.

  \item \(13 - (3 - 7) = 13 - (-4) = 13 + 4 = \mathbf{17}\) % 2.

  \item \(12 - (-2 - 6) = 12 - (-8) = 12 + 8 = \mathbf{20}\) % 3.

  \item \(7 - (-4 + 6) = 7 - 2 = \mathbf{5}\) % 4.

  \item \(1 + (3 - 7) = 1 + (-4) = 1 - 4 = \mathbf{-3}\) % 5.

  \item \(4 + (-3 - 7) = 4 + (-10) = 4 - 10 = \mathbf{-6}\) % 6.

  \item \(0 - (8 - 3) = 0 - 5 = \mathbf{-5}\) % 7.

  \item \(0 - (-4 - 9) = 0 - (-13) = \mathbf{13}\) % 8.

  \item \(-(-2 + 6) = \mathbf{-4}\) % 9.

  \item \(-(4 + 6) = \mathbf{-10}\) % 10.

  \item \(-(5 - 12) = -(-7) = \mathbf{7}\) % 11.

  \item \(-(-3 - 7) - (8 - 5) = -(-10) - 3 = 10 - 3 = \mathbf{7}\) % 12.

  \item \(6 - (8 - 5) - 3 = 6 - 3 - 3 = 3 - 3 = \mathbf{0}\) % 13.

  \item \(-(4 - 7) - (-3 - 2) = -(-3) - (-5) = 3 + 5 = \mathbf{8}\) % 14.

  \item \(-(-3 + 7) + (8 - 5) + 3 = -4 + 3 + 3 = -1 + 3 = \mathbf{2}\) % 15.

  \item \((4 - 8) - (9 - 6) = -4 - 3 = \mathbf{-7}\) % 16.

  \item \(-(-3 + 11) + (8 - 5) = -8 + 3 = \mathbf{-5}\) % 17.

  \item \(-(-2 - 6) - (13 - 7) = -(-8) - 6 = 8 - 6 = \mathbf{2}\) % 18.

  \item \(-(9 + 4) - (-11 + 17) = -13 - 6 = \mathbf{-19}\) % 19.

  \item \(-(9 - 4 + 11 - 17) = -(5 + 11 - 17) = -(16 - 17) = -(-1) = \mathbf{1}\) % 20.

  \item \(7 - (4 + 6) = 7 - 10 = \mathbf{-3}\) % 21.

  \item \(-(2 - (6 - 8)) = -(2 - (-2)) = -(2 + 2) = \mathbf{-4}\) % 22.

\end{enumerate}

\subsubsection{Schreibe übersichtlich und berechne:} % IV.

\begin{enumerate}

  \item \(3 - \big( 7 - (10 - 4) \big) = 3 - (7 - 6) = 3 - 1 = \mathbf{2}\) % 1.

  \item \(14 - \big( 5 - (8 - 3) \big) = 14 - (5 - 5) = 14 - 0 = \mathbf{14}\) % 2.

  \item \(12 - \big( -4 - (-2 - 6) \big) = 12 - \big( -4 - (-8) \big) = 12 - (-4 + 8) = 12 - 4 = \mathbf{8}\) % 3.

  \item \(3 - \big( -(-4 + 6) + 8 \big) = 3 - (-2 + 8) = 3 - 6 = \mathbf{-3}\) % 4.

  \item \(9 - \big( -(3 - 7) - (-5 + 13) \big) = 9 - \big( -(-4) - 8 \big) = 9 - (4 - 8) = 9 - (-4) = 9 + 4 =
    \mathbf{13}\) % 5.

  \item \(- \big( 4 - (6 - 7) \big) + 7 = - \big( 4 - (-1) \big) + 7 = -(4 + 1) + 7 = -5 + 7 = \mathbf{2}\) % 6.

  \item \((-4 - 9) - (8 - 3) = -13 - 5 = \mathbf{-18}\) % 7.

  \item \(-(-5 + 3) + (7 - 18) = -(-2) + (-11) = 2 - 11 = \mathbf{-9}\) % 8.

  \item \(- \big( (-5 + 3) + (7 - 18) \big) = - \big( -2 + (-11) \big) = -(-2 - 11) = -(-13) = \mathbf{13}\) % 9.

  \item \(3 - \Big( 4 - \big( 5 - (-7 + 8) \big) \Big) = 3 - \big( 4 - (5 - 1) \big) = 3 - (4 - 4) = 3 - 0 =
    \mathbf{3}\) % 10.

  \item \((-8 + 15) - \big( 5 - (6 - 18) \big) = 7 - \big( 5 - (-12) \big) = 7 - (5 + 12) = 7 - 17 =
    \mathbf{-10}\) % 11.

  \item \(- \big( (-4 - 9) - (12 - 7) \big) = -(-13 - 5) = -(-18) = \mathbf{18}\) % 12.

  \item \(6 - \big( (8 - 5) - 3 \big) = 6 - (3 - 3) = 6 - 0 = \mathbf{6}\) % 13.

  \item \(1 + \Big( \big( -(4 - 7) + 3 \big) - \big( -1 - 2 \big) \Big) =
    1 + \Big( \big( -(-3) + 3 \big) - \big( -1 - 2 \big) \Big) =
    1 + \big( (3 + 3) - (-3) \big) = 1 + (6 + 3) = 1 + 9 = \mathbf{10}\) % 14.

  \item \(\big( -(-5) \big) + (-5) = 5 - 5 = \mathbf{0}\) % 15.

  \item \(4 - \Big( - \big( -(-1) \big) \Big) = 4 - (-1) = 4 + 1 = \mathbf{5}\) % 16.

  \item \(5 + \big( 4 - (-3 + 11) \big) = 5 + (4 - 8) = 5 + (-4) = 5 - 4 = \mathbf{1}\) % 17.

  \item \(4 - (-3 - 9) - (12 - 25) = 4 - (-12) - (-13) = 4 + 12 + 13 = 16 + 13 = \mathbf{29}\) % 18.

  \item \(5 - \Big( \big( -3 + 2 \big) - \big( 7 + (-9) \big) \Big) = 5 - \big( -1 - (7 - 9) \big) = 5 - \big( -1 - (-2) \big) = 5 - (-1 + 2) = 5 - 1 =
    \mathbf{4}\) % 19.

  \item \(1 - \bigg( 2 - 3 + \Big( 4 - \big( 5 - 6 + (7 - 8) \big) \Big) \bigg) = 1 - \bigg( -1 + \Big( 4 - \big( -1 + (-1) \big) \Big) \bigg) =
    1 - \Big( -1 + \big( 4 - (-1 - 1) \big) \Big) = 1 - \Big( -1 + \big( 4 - (-2) \big) \Big) = 1 - (-1 + 6) = 1 - 5 = 1 + 3 = \mathbf{-4}\) % 20.

\end{enumerate}

\subsection{Aufgaben zu Kapitel 2.9 bis 2.10 \\
  - Multiplikation mit negativen Zahlen \\
  - Division zweier Zahlen}

\subsubsection{Berechne:} % I.

\begin{enumerate}[label=\alph*)]

  \item \(3 \cdot (-5) = \mathbf{-15}\) % a)

  \item \((-4) \cdot 8 = \mathbf{-32}\) % b)

  \item \((-4) \cdot (-4) = \mathbf{16}\) % c)

  \item \(-(-2) \cdot (-3) = 2 \cdot (-3) \mathbf{-6}\) % d)

  \item \((-5) \cdot (-6) = \mathbf{30}\) % e)

  \item \(8 \cdot (-4) + (-2) \cdot (-3) = -32 + 6 = \mathbf{-26}\) % f)

  \item \((-3) \cdot 5 + (-7) \cdot (-4) = -15 + 28 = \mathbf{13}\) % g)

  \item \((-3) \cdot (-2) + 2 \cdot (-1) + 5 = 6 + (-2) + 5 = 4 + 5 = \mathbf{9}\) % h)

  \item \((-5) \cdot 3 + 9 + (-4) \cdot (-2) = -15 + 9 + 8 = -6 + 8 = \mathbf{2}\) % i)

  \item \((-3) \cdot 8 - (-7) \cdot (-5) = -24 - 35 = \mathbf{-59}\) % j)

  \item \((-1) \cdot (-1) = \mathbf{1}\) % k)

  \item \((-1) \cdot 1 = \mathbf{-1}\) % l)

  \item \(1 \cdot (-1) = \mathbf{-1}\) % m)

  \item \((-1) \cdot 8 + (-1) \cdot 5 = -8 - 5 = \mathbf{-13}\) % n)

  \item \((-1) \cdot 9 + (-1) \cdot 4 + 7 + (-1) \cdot 3 = -9 - 4 + 7 - 3 =
    -13 + 7 - 3 = -6 - 3 = \mathbf{-9}\) % o)

  \end{enumerate}

\subsubsection{Berechne:} % II.

\begin{enumerate}

  \item \(54 \colon (-3) = \mathbf{-18}\) % 1.

  \item \((-48) \colon 8 = \mathbf{-6}\) % 2.

  \item \((-56) \colon (-7) = \mathbf{8}\) % 3.

  \item \((-144) \colon 6 = \mathbf{-24}\) % 4.

  \item \((-54) \colon (-6) = \mathbf{9}\) % 5.

  \item \(\cfrac{-3}{4} = \mathbf{-\cfrac{3}{4}}\) % 6.

  \item \(\cfrac{10}{-16} = \mathbf{-\cfrac{5}{8}}\) % 7.

  \item \(\cfrac{-10}{-12} = \mathbf{\cfrac{5}{6}}\) % 8.

  \item \(\cfrac{-2 + 3 - 7 + 4}{5 - 8 + 2 - 6 + 1} = \cfrac{1 - 3}{-3 - 4 + 1} =
    \cfrac{-2}{-6} = \mathbf{\cfrac{1}{3}}\) % 8.

  \item \(\cfrac{3 - 8 + 1 - 11 + 5}{5 - 8 + 2 - 6 + 1 - (-10)} = \cfrac{-5 - 10 + 5}{-3 - 4 + 11} =
    \cfrac{-15 + 5}{-7 + 11} = \cfrac{-10}{4} = \mathbf{-\cfrac{5}{2}}\) % 9.

  \item \(\cfrac{1 + 4 - 1 + 4 - 2 + 1 - 3 + 5}{1 - 7 + 3 - 20 + 50 - 8} =
    \cfrac{5 + 3 - 1 + 2}{-6 - 17 + 42} = \cfrac{8 - 3}{-23 + 42} = \mathbf{\cfrac{5}{19}}\) % 10.

  \item \(\cfrac{-3 - 1 + 4 + 1 - 5 + 9 - 2 - 7}{2 + 2 - 3 - 6 - 0 + 6 - 8} =
    \cfrac{-4 + 5 + 4 - 9}{4 - 9 - 2} = \cfrac{1 - 5}{-5 - 2} = \cfrac{-4}{-7} = \mathbf{\cfrac{4}{7}}\) % 11.

  \item \(\cfrac{-4 - 7 - 7 + 1 - 2 + 1 - 3}{6 - 9 - 3 + 1 - 4 + 7 - 2} =
    \cfrac{-11 - 6 - 1 - 2}{-3 - 2 + 3} = \cfrac{-17 - 3}{-5 + 3} = \cfrac{-21}{-2} = \mathbf{\cfrac{21}{2}}\) % 12.

  \item \(\cfrac{3 - 8 + 1 - 11 + 5}{5 - 8 + 2 - 6 + 1 - (-10)} = \cfrac{-5 - 10 + 5}{-3 - 4 + 11} =
    \cfrac{-15 + 5}{-7 + 11} = \cfrac{-10}{4} = \mathbf{-\cfrac{5}{2}}\) % 13.

  \item \(\cfrac{2 - 7 + 1 - 8 + 2 - 8 + 1 - 8}{3 - 1 - 6 - 2 - 2 + 7 + 7 - 7} =
    \cfrac{-5 - 7 - 6 - 7}{2 - 8 + 5} = \cfrac{-12 - 13}{-6 + 5} = \cfrac{-25}{-1} = \mathbf{25}\) % 14.

  \item \(\cfrac{-3}{4} + \cfrac{2}{-5} = -\cfrac{3}{4} - \cfrac{2}{5} = -\cfrac{15}{20} - \cfrac{8}{20} =
    \mathbf{-\cfrac{23}{20}}\) % 15.

  \item \(\cfrac{-4}{-5} + \cfrac{2}{-8} = \cfrac{4}{5} - \cfrac{1}{4} = \cfrac{16}{20} - \cfrac{5}{20} =
    \mathbf{\cfrac{11}{20}}\) % 16.

  \item \(\cfrac{4}{-7} - \cfrac{-3}{5} = -\cfrac{4}{7} + \cfrac{3}{5} = -\cfrac{20}{35} + \cfrac{21}{35} =
    \mathbf{\cfrac{1}{35}}\) % 17.

  \item \(\cfrac{-1}{9} - \cfrac{1}{-11} = -\cfrac{1}{9} + \cfrac{1}{11} = -\cfrac{11}{99} + \cfrac{9}{99} =
    \mathbf{-\cfrac{2}{99}}\) % 18.

  \item \(\cfrac{5}{-13} + \cfrac{4}{-5} = -\cfrac{5}{13} - \cfrac{4}{5} = -\cfrac{25}{65} - \cfrac{52}{65} =
    \mathbf{-\cfrac{77}{65}}\) % 19.

\end{enumerate}

\subsection{Aufgaben zu Kapitel 2.11 \\
  - Potenzen}

\subsubsection{Berechne:} % I.

\begin{enumerate}[label=\alph*)]

  \item \(2^{-5} = \cfrac{1}{2^5} = \cfrac{1}{2^3} \cdot \cfrac{1}{2^2} = \cfrac{1}{8} \cdot \cfrac{1}{4} =
    \mathbf{\cfrac{1}{32}}\) % a)

  \item \((-3)^{-4} = -\cfrac{1}{3^4} = -\cfrac{1}{3^2} \cdot \cfrac{1}{3^2} =
    -\cfrac{1}{9} \cdot \cfrac{1}{9} = -\cfrac{1}{81}\) % b)

  \item \((-3)^{-3} = -\cfrac{1}{3^3} = -\cfrac{1}{3^2} \cdot \cfrac{1}{3^1} =
    -\cfrac{1}{9} \cdot \cfrac{1}{3} = \mathbf{-\cfrac{1}{27}}\) % c)

  \item \(5^{-4} = \cfrac{1}{5^4} = \cfrac{1}{5^2} \cdot \cfrac{1}{5^2} =
    \cfrac{1}{25} \cdot \cfrac{1}{25} = \cfrac{1}{125} \cdot \cfrac{1}{5} = \mathbf{\cfrac{1}{625}}\) % d)

  \item \(2^{-10} = \cfrac{1}{2^{10}} = \cfrac{1}{2^5} \cdot \cfrac{1}{2^5} =
    \cfrac{1}{32} \cdot \cfrac{1}{32} = \cfrac{1}{64} \cdot \cfrac{1}{16} = \cfrac{1}{128} \cdot \cfrac{1}{8} =
    \cfrac{1}{256} \cdot \cfrac{1}{4} = \cfrac{1}{512} \cdot \cfrac{1}{2} = \mathbf{\cfrac{1}{1024}}\) % e)

  \item \(10^{-3} = \cfrac{1}{10^{3}} = \mathbf{\cfrac{1}{1000}}\) % f)

  \item \(2^9 \cdot 2^8 \cdot 2^{-14} = 2^{9 + 8 - 14} = 2^{17 - 14} = 2^3 = \mathbf{8}\) % g)

  \item \(2^7 \cdot 3^7 \cdot 2^{-5} \cdot 3^{-5} = 2^{7 - 5} \cdot 3^{7 - 5} =
    2^2 \cdot 3^2 = 4 \cdot 9 = \mathbf{36}\) % h)

  \item \(2^{93} \cdot 2^{27} \cdot 2^{-42} \cdot 2^{-69} = 2^{93 + 27 - 42 - 69} =
    2^{120 - 111} = 2^9 = 2^5 \cdot 2^4 = 32 \cdot 16 = 64 \cdot 8 = 128 \cdot 4 = 256 \cdot 2 =
    \mathbf{512}\) % i)

  \item \(2^{23} \cdot 3^2 \cdot 2^{-3} \cdot 2^{-16} = 2^{23 - 3 - 16} \cdot 3^2 =
    2^4 \cdot 3^2 = 16 \cdot 9 = \mathbf{144}\) % j)

  \item \(2^5 + 2^5 \cdot 3^5 \cdot 3^{-5} = 32 + 32 \cdot 3^{5 - 5} = 32 + 32 \cdot 3^0 =
    32 + 32 \cdot 1 = \mathbf{64}\) % k)

  \item \(2^{-3} \cdot 3^2 = \cfrac{3^2}{2^3} = \mathbf{\cfrac{9}{8}}\) % l)

\end{enumerate}

\subsubsection{Berechne:} % II.

\begin{enumerate}[label=\alph*)]

  \item \(\cfrac{3^5 \cdot 3^8}{3^2 \cdot 3^4} = \cfrac{3^{5 + 8}}{3^{2 + 4}} =
    \cfrac{3^{13}}{3^6} = 3^{13 - 6} = 3^7 = 3^4 \cdot 3^3 = 81 \cdot 27 = 243 \cdot 9 = \mathbf{2.187}\) % a)

  \item \(\cfrac{2^7 \cdot 2^4}{2^5} = \cfrac{2^{7 + 4}}{2^5} = \cfrac{2^{11}}{2^5} =
    2^{11 - 5} = 2^6 = 2^3 \cdot 2^3 = 8 \cdot 8 = \mathbf{64}\) % b)

  \item \(\cfrac{2^5 \cdot 3^5}{2^3 \cdot 3^2} = 2^{5 - 3} \cdot 3^{5 - 2} =
    2^2 \cdot 3^3 = 4 \cdot 27 = \mathbf{108}\) % c)

  \item \(\cfrac{(2^3)^2}{2^2 \cdot 2^3} = \cfrac{2^{3 \cdot 2}}{2^{2 + 3}} =
    \cfrac{2^6}{2^5} = 2^{6 - 5} = 2^1 = \mathbf{2}\) % d)

  \item \(\cfrac{(2^2)^2 \cdot (3^2)^2}{2^2 \cdot 3^2} = 2^{2 \cdot 2 - 2} \cdot 3^{2 \cdot 2 - 2} =
    2^{4 - 2} \cdot 3^{4 - 2} = 2^2 \cdot 3^2 = 4 \cdot 9 = \mathbf{36}\) % e)

  \item \(\cfrac{2^4 \cdot 3^4 \cdot 2^5}{3^3 \cdot (3^2)^2} = 2^{4 + 5} \cdot 3^{4 - 3 - 2 \cdot 2} =
    2^9 \cdot 3^{1 - 4} = 2^9 \cdot 3^{-3} = \cfrac{2^9}{3^3} = \mathbf{\cfrac{512}{27}}\) % f)

\end{enumerate}

\subsection{Aufgaben zu Kapitel 2.12 \\
  - Wissenschaftliche Notation von Zahlen}

\subsubsection{Berechne:} % I.

\begin{enumerate}[label=\alph*)]

  \item \(3,517 \cdot 10^3 = \mathbf{3.517}\) % a)

  \item \(1,4142135 \cdot 10^{-3} = \mathbf{0,001.414.213.5}\) % b)

  \item \(1,73205 \cdot 10^{-1} = \mathbf{0,173.205}\) % c)

  \item \(3,14159 \cdot 10^{12} = \mathbf{3.141.590.000.000}\) % d)

  \item \(2,7182818 \cdot 10^{-7} = \mathbf{0,000.000.271.828.18}\) % e)

  \item \(3,1622776 \cdot 10^{-1} = \mathbf{0,316.227.76}\) % f)

  \item \(8660,254 \cdot 10^{-4} = \mathbf{0,866.025.4}\) % g)

  \item \(7071067,8 \cdot 10^{-8} = \mathbf{0,070.710.678}\) % h)

  \item \(0,30102999 \cdot 10^{-6} = \mathbf{0,000.000.301.029.99}\) % i)

  \item \(693,147180 \cdot 10^7 = \mathbf{6.931.471.800}\) % j)

\end{enumerate}

\subsubsection{Schreibe in wissenschaftlicher Notation:} % II.

\begin{enumerate}[label=\alph*)]

  \item \(7.325 = \mathbf{7,325 \cdot 10^3}\) % a)

  \item \(1.358.732 = \mathbf{1,358.732 \cdot 10^6}\) % b)

  \item \(972.144 = \mathbf{9,721.44 \cdot 10^5}\) % c)

  \item \(0,141.537 = \mathbf{1,415.37 \cdot 10^{-1}}\) % d)

  \item \(0,053.428.9 = \mathbf{5,342.89 \cdot 10^{-2}}\) % e)

  \item \(0,000.024.698 = \mathbf{2,469.8 \cdot 10^{-5}}\) % f)

  \item \(14.031.879 = \mathbf{1,403.187.9 \cdot 10^7}\) % g)

  \item \(0,001.710.927 = \mathbf{1,710.927 \cdot 10^{-3}}\) % h)

  \item \(0,000.000.000.08 = \mathbf{8,0 \cdot 10^{-11}}\) % i)

  \item \(432.177.605 = \mathbf{4,321.776.05 \cdot 10^8}\) % j)

\end{enumerate}

\section{Buchstabenrechnung}

\subsection{Aufagaben zu Kapitel 3.1 \\
  - Addition und Subtraktion von Buchstaben}

\subsubsection{Berechne:} % I.

\begin{enumerate}[label=\alph*)]

  \item \(8a - 5a = \mathbf{3a}\) % a)

  \item \(-5a + 8a = \mathbf{3a}\) % b)

  \item \(-7a + 3a = \mathbf{-4a}\) % c)

  \item \(13a - 19a = \mathbf{-6a}\) % d)

  \item \(-21a + 8a = \mathbf{-13a}\) % e)

  \item \(-8b - 5b = \mathbf{-13b}\) % f)

  \item \(8b - 5b + 7b - 20b = 3b - 13b = \mathbf{-10b}\) % g)

  \item \(-9c - 4c + 11c - 17c = -13c - 6c = \mathbf{-19c}\) % h)

  \item \(-(9a + 4a) - (-11a + 17a) = -13a - 6a = \mathbf{-19a}\) % i)

  \item \(-(9a - 4a + 11a - 17a) = -(5a - 6a) = -(-a) = \mathbf{a}\) % j)

  \item \(4x - 9x = \mathbf{-5x}\) % k)

  \item \(-5x + 18x = \mathbf{13x}\) % l)

  \item \(12x - 7x = \mathbf{5x}\) % m)

  \item \(8x - 5y - 3x + 2y = \mathbf{5x - 3y}\) % n)

  \item \(x - 6y + 4y - 3x = \mathbf{-2x - 2y}\) % o)

  \item \(8a - 5b + 2c - 3a + 7b - 4c = \mathbf{5a + 2b - 2c}\) % p)

\end{enumerate}

\subsubsection{Berechne:} % II.

\begin{enumerate}[label=\alph*)]

  \item \((4a + 5b) + (2a + 3b) = 4a + 5b + 2a + 3b = \mathbf{6a + 8b}\) % a)

  \item \((4a + 5b) + (2a - 3b) = 4a + 5b + 2a - 3b = \mathbf{6a + 2b}\) % b)

  \item \((4a + 5b) - (2a - 3b) = 4a + 5b - 2a + 3b = \mathbf{2a + 8b}\) % c)

  \item \((4a + 5b) - (-2a + 3b) = 4a + 5b + 2a - 3b = \mathbf{6a + 2b}\) % d)

  \item \((4a + 5b) - (-2a - 3b) = 4a + 5b + 2a + 3b = \mathbf{6a + 8b}\) % e)

  \item \(-(9a - 4b) - (-11a  + 17b) = -9a + 4b + 11a - 17b = \mathbf{2a - 13b}\) % f)

  \item \(7a - (4b + 6a) = 7a - 4b - 6a = \mathbf{a - 4b}\) % g)

  \item \(- \big( 2a - (6a - 8b) \big) = -(2a - 6a + 8b) = -2a + 6a - 8b = \mathbf{4a - 8b}\) % h)

  \item \(5a - \Big( \big( -3b + 2a \big) - \big(7a + (-9b) \big) \Big) = 5a - \big( -3b + 2a - 7a - (-9b) \big) =
    5a - (6b - 5a) = 5a - 6b + 5a = \mathbf{10a - 6b}\) % i)

  \item \(x - \bigg[ 2y - 3x + \Big( 4x - \big( 5y - 6x + (7x - 8y) \big) \Big) \bigg] =
    x - 2y + 3x - \big( 4x - 5y + 6x - (7x - 8y) \big) = 4x - 2y - (10x - 5y - 7x + 8y) = 4x - 2y - (3x + 3y) =
    4x - 2y - 3x - 3y = \mathbf{x - 5y}\) % j)

\end{enumerate}

\subsubsection{Berechne:} % III.

\begin{enumerate}[label=\alph*)]

  \item \(6c - 3c = \mathbf{3c}\) % a)

  \item \(7d - 11d = \mathbf{-4d}\) % b)

  \item \(-5x + 8x = \mathbf{3x}\) % c)

  \item \(-4y - 7y = \mathbf{-11y}\) % d)

  \item \(5c - 2b + 4a + 1 - 3b + 2a - 7c - 3 = \mathbf{6a - 5b - 2c -2}\) % e)

  \item \(-8y + 2z - x - 4z + 2y + 5x = \mathbf{4x - 6y - 2z}\) % f)

  \item \(4w - 2v + 8u - 5w - 11u = \mathbf{-3u - 2v - w}\) % g)

  \item \(3n - 2m + 2n - 4n - 9n = (3 + 2 - 4 - 9)n - 2m = (5 - 13)n - 2m = \mathbf{-2m - 8n}\) % h)

\end{enumerate}

\subsubsection{Berechne:} % IV.

\begin{enumerate}[label=\alph*)]

  \item \(2a - 3b + (4b - a) - (7a - 2b - 3b) = 2a - 3b + 4b - a - (7a - 5b) =
    2a + b - a - 7a + 5b = (2 - 1 - 7)a + (-1 + 5)b = \mathbf{-6a + 4b}\) % a)

  \item \(2b - 3c + 4a - (-6c + 2a - 2c + 3b - 7a) = 2b - 3c + 4a - (-8c - 5a + 3b) =
    2b - 3c + 4a + 8c + 5a - 3b = \mathbf{9a - 1b + 5c}\) % b)

  \item \(-(-2x + 7y - 3x) + (5y - 6x + 2x - 3y) - 3x + 9y = -(-5x + 7y) + (2y - 4x) - 3x + 9y =
    5x - 7y + 2y - 4x - 3x + 9y = (5 - 4 - 3)x + (-7 + 2 + 9)y = \mathbf{-2x + 4y}\) % c)

  \item \(-3y - 4x - (2y - x + 3y) = -3y - 4x - (5y - x) = -3y - 4x - 5y + x = (-4 + 1)x + (-3 - 5)y =
    \mathbf{-3x - 8y}\) % d)

  \item \(6 - 3b + 2a - \big( 2b - 3a - (4a - 3b) \big) = 6 - 3b + 2a - (2b - 3a - 4a + 3b) =
    6 - 3b + 2a - (5b - 7a) = 6 - 3b + 2a - 5b + 7a = (2 + 7)a + (-3 - 5)b + 6 = \mathbf{9a - 8b + 6}\) % e)

  \item \(-7x + 3 - (3x + 1 - 7x + 3) = -7x + 3 - (-4x + 4) = -7x + 3 + 4x - 4 = \mathbf{-3x - 1}\) % f)

  \item \(2 - (3v - 5w) + (2w - 4 - v + 5w + 1) = 2 - 3v + 5w + (7w - 3 - v) = 2 - 3v + 5w + (7w - 3 - v) =
    2 - 3v + 5w + 7w - 3 - v = (-3 - 1)v + (5 + 7)w - 1 = \mathbf{-4v + 12w - 1}\) % g)

  \item \(\big( 2 - (u - 3v) + 5u \big) - \big( 3v + 1 - (-2v + 3u) \big) =
    (2 - u + 3v + 5u) - (3v + 1 + 2v - 3u) = (2 + 4u + 3v) - (5v + 1 - 3u) = 2 + 4u + 3v - 5v - 1 + 3u =
    (4 + 3)u + (3 - 5)v + 1 = \mathbf{7u - 2v + 1}\) % h)

\end{enumerate}

\subsection{Aufgaben zu Kapitel 3.2 bis 3.4 \\
  - Reine Multiplikation von Buchstaben \\
  - Multiplikation von Summen und Differenzen \\
  - Gemischte Beispiele}

\subsubsection{Multipliziere aus:} % I.

\begin{enumerate}[label=\alph*)]

  \item \(x \cdot (x + 5) = \mathbf{x^2 + 5x}\) % a)

  \item \((x + 3) \cdot x = \mathbf{x^2 + 3x}\) % b)

  \item \(x \cdot (x - 7) = \mathbf{x^2 - 7x}\) % c)

  \item \(a \cdot (8 - a) = \mathbf{-a^2 + 8a}\) % d)

  \item \(-3 \cdot (x - 12) = \mathbf{-3x + 36}\) % e)

  \item \(a \cdot (-x - 5) = \mathbf{-ax - 5a}\) % f)

  \item \(-a \cdot (-x + 5) = \mathbf{ax - 5a}\) % g)

  \item \(x \cdot (8x - 5z - 3 + 2y) = \mathbf{8x^2 + 2xy - 5xz - 3x}\) % h)

  \item \((a - 2b + 3c) \cdot a = \mathbf{a^2 - 2ab + 3ac}\) % i)

  \item \((a - b + c) \cdot (-a) = \mathbf{-a^2 + ab - ac}\) % j)

  \item \((a + 3) \cdot (b + 4) = \mathbf{ab + 4a + 3b + 12}\) % k)

  \item \((a + 2) \cdot (-b + 5) = \mathbf{-ab + 5a - 2b + 10}\) % l)

  \item \((x + 6) \cdot (x + 3) = x^2 + 3x + 6x + 18 = \mathbf{x^2 + 9x + 18}\) % m)

  \item \((a + 1) \cdot (b + 2) = \mathbf{ab + 2a + b + 2}\) % n)

  \item \((2x + 3) \cdot (5x - 7) = 10x^2 - 14x + 15x - 21 = \mathbf{10x^2 + x - 21}\) % o)

  \item \((2x + 3) \cdot (3x - 5) = 6x^2 - 10x + 9x - 15 = \mathbf{6x^2 - x - 15} \) % p)

\end{enumerate}

\subsubsection{Multipliziere aus:} % II.

\begin{enumerate}[label=\alph*)]

  \item \((x + 1) \cdot (x + 2) = x^2 + 2x + x + 2 = \mathbf{x^2 + 3x + 2}\) % a)

  \item \((x + 4) \cdot (x + 5) = x^2 + 5x + 4x + 20 = \mathbf{x^2 + 9x + 20}\) % b)

  \item \((x + 2) \cdot (x - 3) = x^2 - 3x + 2x - 6 = \mathbf{x^2 - x - 6}\) % c)

  \item \((x - 3) \cdot (x + 4) = x^2 + 4x - 3x - 12 = \mathbf{x^2 + x - 12}\) % d)

  \item \((x - 5) \cdot (x - 1) = x^2 - x - 5x + 5 = \mathbf{x^2 - 6x + 5}\) % e)

  \item \((-x + 3) \cdot (x - 4) = -x^2 + 4x + 3x - 12 = \mathbf{-x^2 + 7x - 12}\) % f)

  \item \((-x - 1) \cdot (-x + 2) = x^2 - 2x + x - 2 = \mathbf{x^2 - x - 2}\) % g)

  \item \((-x - 3) \cdot (-x - 7) = x^2 + 7x + 3x + 21 = \mathbf{x^2 + 10x + 21}\) % h)

  \item \((a - 2) \cdot (b + 5) = \mathbf{ab + 5a - 2b - 10}\) % i)

  \item \((u + 1) \cdot (-v + 6) = \mathbf{-uv + 6u - v + 6}\) % j)

  \item \((x - 4) \cdot (b - 3) = \mathbf{bx -4b - 3x + 12}\) % k)

  \item \((z + 2) \cdot (w - 1) = \mathbf{wz + 2w - z - 2}\) % l)

  \item \((x - 7) \cdot (y + 6) = \mathbf{xy + 6x - 7y - 42}\) % m)

  \item \((k - 2) \cdot (k - 3) = k^2 - 3k - 2k + 6 = \mathbf{k^2 - 5k + 6}\) % n)

  \item \((m + 4) \cdot (n - 4) = \mathbf{mn - 4m + 4n - 16}\) % o)

  \item \((a + 3) \cdot (a + 3) = \mathbf{a^2 + 6a + 9}\) % p)

\end{enumerate}

\subsubsection{Multipliziere aus:} % III.

\begin{enumerate}[label=\alph*)]

  \item \((3x + 2) \cdot (4x + 5) = 12x^2 + 15x + 8x + 10 = \mathbf{12x^2 + 23x + 10}\) % a)

  \item \((4x + 3) \cdot (2x + 1) = 8x^2 + 4x + 6x + 3 = \mathbf{8x^2 + 10x + 3}\) % b)

  \item \((5x + 2) \cdot (3x - 4) = 15x^2 - 20x + 6x - 8 = \mathbf{15x^2 - 14x - 8}\) % c)

  \item \((2x - 3) \cdot (y + 4) = \mathbf{2xy + 4x - 3y - 12}\) % d)

  \item \((3y - 5) \cdot (2x - 6) = \mathbf{6xy -10x - 18y + 30}\) % e)

  \item \((-2x + 3y) \cdot (5x - 4y) = -10x^2 + 8xy + 15xy -12y^2 = \mathbf{-10x^2 + 23xy - 12^2}\) % f)

  \item \((-x - y) \cdot (-2x + y) = 2x^2 - xy + 2xy - y^2 = \mathbf{2x^2 + xy - y^2}\) % g)

  \item \((-a - 2b) \cdot (-a - 6b) = a^2 + 6ab + 2ab + 12b^2 = \mathbf{a^2 + 8ab + 12b^2}\) % h)

  \item \((3a - 2b) \cdot (5a + b) = 15a^2 + 3ab - 10ab - 2b^2 = \mathbf{15a^2 - 7ab - 2b^2}\) % i)

  \item \((u + v) \cdot (-u + 4v) = -u^2 + 4uv - uv + 4v^2 = \mathbf{-u^2 + 3uv + 4v^2}\) % j)

  \item \((2u - 4w) \cdot (3v - 5w) = \mathbf{6uv - 10uw - 12vw + 20w^2}\) % k)

  \item \((w + 2) \cdot (z - 1) = \mathbf{wz - w + 2z - 2}\) % l)

  \item \((5x - 6y) \cdot (2x + 3y) = 10x^2 + 15xy -12xy - 18y^2 = \mathbf{10x^2 + 3xy - 18y^2}\) % m)

  \item \((a - 2b + 3c) \cdot (2a - b + 5) = 2a^2 - ab + 5a - 4ab + 2b^2 - 10b + 6ac - 3bc + 15c =
  \mathbf{2a^2 - 5ab + 6ac + 2b^2 - 3bc + 5a - 10b + 15c}\) % n)

  \item \((2m + 3n) \cdot (3m - 2n) = 6m^2 - 4mn + 9mn - 6n^2 = \mathbf{6m^2 + 5mn - 6n^2}\) % o)

  \item \((2a + 3) \cdot (2a - 3) = \mathbf{4a^2 - 9}\) % p)

\end{enumerate}

\subsubsection{Multipliziere aus:} % IV.

\begin{enumerate}[label=\alph*)]

  \item \(\left( \cfrac{3}{2}x + \cfrac{1}{2}y \right) \cdot \left( \cfrac{5}{2}x - \cfrac{1}{2}y \right) =
    \cfrac{15}{4}x^2 - \cfrac{3}{4}xy + \cfrac{5}{4}xy - \cfrac{1}{4}y^2 =
    \mathbf{\cfrac{15}{4}x^2 + \cfrac{1}{2}xy - \cfrac{1}{4}y^2}\) % a)

  \item \(\left( -\cfrac{2}{5}x + \cfrac{1}{3}y \right) \cdot \left( \cfrac{3}{4}x - \cfrac{5}{6}y \right) =
    -\cfrac{6}{20}x^2 + \cfrac{10}{30}xy + \cfrac{3}{12}xy - \cfrac{5}{18} =
    -\cfrac{3}{10}x^2 + \cfrac{1}{3}xy + \cfrac{1}{4}xy - \cfrac{5}{18} =
    \mathbf{-\cfrac{3}{10}x^2 + \cfrac{7}{12}xy - \cfrac{5}{18}y^2}\) % b)

  \item \(x \cdot \left( 1 + \cfrac{1}{2}x \cdot \left( 1 + \cfrac{1}{3}x \right) \right) =
    x \cdot \left( 1 + \cfrac{1}{2}x + \cfrac{1}{6}x^2 \right) = \mathbf{\cfrac{1}{6}x^3 + \cfrac{1}{2}x^2 + x}\) % c)

  \item \((0,2x - 1,3y) \cdot (0,5x + 4y) = 0,1x^2 + 0,8xy - 0,65xy - 5,2y^2 = \mathbf{0,1x^2 - 0,15xy - 5,2y^2}\) % d)

  \item \((x - 1) \cdot (1 + x + x^2 + x^3) = x + x^2 + x^3 + x^4 - 1 - x - x^2 - x^3 = \mathbf{x^4 - 1}\) % e)

  \item \((x + 1) \cdot (1 - x + x^2 - x^3) = x - x^2 + x^3 - x^4 + 1 - x + x^2 - x^3 = \mathbf{-x^4 + 1}\) % f)

  \item \((4x^2 - 9y^2) \cdot (4x^2 + 9y^2) = (4x^2)^2 - (9y^2)^2 = \mathbf{16x^4 - 81y^4}\) % g)

  \item \((u^2 - 3v) \cdot (-2u - 5v^2) = \mathbf{-5u^2v^2 - 2u^3 + 15v^3 + 6uv}\) % h)

  \item \((2a - 3c) \cdot (4b + 2c) = \mathbf{8ab + 4ac - 12bc - 6c^2}\) % i)

  \item \((5u - 3v + 2w) \cdot (-u + 2v + 3w) = -5u^2 + 10uv + 15uw + 3uv - 6v^2 - 9vw - 2uw + 4vw + 6w^2 =
    \mathbf{-5u^2 + 13uv + 13uw - 6v^2 - 5vw + 6w^2}\) % j)

  \item \((4x - 2y + 3) \cdot (5x - 3y + 2) = 20x^2 - 12xy + 8x - 10xy + 6y^2 - 4y + 15x - 9y + 6 =
    \mathbf{20x^2 - 22xy + 6y^2 + 23x - 13y + 6}\) % k)

  \item \((x^2 - 3x + 2) \cdot (3y - 4) = \mathbf{3x^2y - 4x^2 - 9xy + 12x + 6y - 8}\) % l)

  \item \((3x - 6y + 2z) \cdot (5x + 3y) = 15x^2 + 9xy - 30xy - 18y^2 + 10xz + 6yz =
    \mathbf{15x^2 - 21xy + 10xz - 18y^2 + 6yz}\) % m)

  \item \((a - 2b + 3c) \cdot (a - 2b + 3c) = a^2 - 4ab + 6ac + 4b^2 - 12bc + 9c^2 =
    \mathbf{a^2 - 4ab + 6ac + 4b^2 - 12bc + 9c^2}\) % n)

  \item \((3u + 2v + 1) \cdot (3u + 2v - 1) = 9u^2 + 6uv - 3u + 6uv + 4v^2 - 2v + 3u + 2v - 1 =
    \mathbf{9u^2 + 12uv + 4v^2 - 1}\) % o)

  \item \((4a - 3) \cdot (4a - 3) = (4a)^2 - 2 \cdot 3 \cdot 4a + 3^2 = \mathbf{16a^2 - 24a + 9}\) % p)

\end{enumerate}

\subsubsection{Multipliziere aus:} % V.

\begin{enumerate}[label=\alph*)]

  \item \(3x \cdot (2x + 5) \cdot x = 3x^2 \cdot (2x + 5) = \mathbf{6x^3 + 15 x^2}\) % a)

  \item \((3x + 5) \cdot (2x + 3) \cdot x = (6x^2 + 9x + 10x + 15) \cdot x = (6x^2 + 19x + 15) \cdot x =
    \mathbf{6x^3 + 19x^2 + 15x}\) % b)

  \item \((x + 5) \cdot (x + 3) \cdot (x - 7) = (x^2 + 8x + 15) \cdot (x - 7) = x^3 - 7x^2 + 8x^2 - 56x + 15x - 105 =
    \mathbf{x^3 + x^2 - 41x - 105}\) % c)

  \item \((x + y + 5) \cdot (x + 2y + 3) = x^2 + 2xy + 3x + xy + 2y^2 + 3y + 5x + 10y + 15 =
    \mathbf{x^2 + 3xy + 2y^2 + 8x + 13y + 15}\) % d)

  \item \((2x - 3y + 5) \cdot (5x + y - 1) = 10x^2 + 2xy - 2x - 15xy - 3y^2 + 3y + 25x + 5y - 5 =
    \mathbf{10x^2 - 13xy - 3y^2 + 23x + 8y - 5}\) % e)

  \item \((x - 1) \cdot (1 + x + x^2) = x + x^2 + x^3 - 1 - x - x^2 = \mathbf{x^3 - 1}\) % f)

  \item \((x - 1) \cdot (x + 2) \cdot (x - 3) = (x^2 + 2x - x - 2) \cdot (x - 3) = (x^2 + x - 2) \cdot (x - 3) =
    x^3 - 3x^2 + x^2 - 3x - 2x + 6 = \mathbf{x^3 - 2x^2 - 5x + 6}\) % g)

  \item \((3x + 1) \cdot (5x - 2) \cdot (2x - 3) = (15x^2 - 6x + 5x - 2) \cdot (2x - 3) =
    (15x^2 - x - 2) \cdot (2x - 3) = 30x^3 - 45x^2 - 2x^2 + 3x - 4x + 6 = \mathbf{30x^3 - 47x^2 - x + 6}\) % h)

  \item \((x^2 - x) \cdot (1 + x + x^2) = x^2 + x^3 + x^4 - x - x^2 - x^3 = \mathbf{x^4 - x}\) % i)

  \item \((x^2 - 1) \cdot (1 + x + x^2 + x^3) = x^2 + x^3 + x^4 + x^5 - 1 - x - x^2 - x^3 =
    \mathbf{x^5 + x^4 - x - 1}\) % j)

  \item \((a + 3) \cdot (b + 4) \cdot (c + 5) = (ab + 4a + 3b + 12) \cdot (c + 5) =
    \mathbf{abc + 5ab + 4ac + 3bc + 20a + 15b + 12c + 60}\) % k)

  \item \((a + b + 2) \cdot (-b + 3 + c) = \mathbf{-ab + ac - b^2 + bc + 3a + b + 2c + 6}\) % l)

  \item \((2a - 3b + 5c) \cdot (a + b - 3 + c) = 2a^2 + 2ab - 6a + 2ac - 3ab - 3b^2 + 9b - 3bc + 5ac + 5bc - 15c + 5c^2 =
    \mathbf{2a^2 - ab + 7ac - 3b^2 + 2bc + 5c^2 - 6a + 9b - 15c}\) % m)

  \item \((a + 1) \cdot (b + 2) \cdot (c - 3) = (ab + 2a + b + 2) \cdot (c - 3) =
    \mathbf{abc - 3ab + 2ac + bc - 6a - 3b + 2c - 6}\) % n)

  \item \((x + 1) \cdot (x + 1) \cdot (x + 1) \cdot (x + 1) = {(x + 1)}^2 \cdot {(x + 1)}^2 =
    (x^2 + 2x + 1) \cdot (x^2 + 2x + 1) = x^4 + 2x^3 + x^2 + 2x^3 + 4x^2 + 2x + x^2 + 2x + 1 =
    \mathbf{x^4 + 4x^3 + 6x^2 + 4x + 1}\) % o)

\end{enumerate}

\subsubsection{Multipliziere aus:} % VI.

\begin{enumerate}[label=\alph*)]

  \item \({(x + 3)}^2 = \mathbf{x^2 + 6x + 9}\) % a)

  \item \((x - 1) \cdot (x + 2) + (x + 3) \cdot (x + 1) = x^2 + x - 2 + x^2 + 4x + 3 =
    \mathbf{2x^2 + 5x + 1}\) % b)

  \item \({(x + 5)}^2 + 3 \cdot (2x + 3) = x^2 + 10x + 25 + 6x + 9 = \mathbf{x^2 + 16x + 34}\) % c)

  \item \(2 \cdot (x + 5) \cdot (x + 3) - (x - 7) \cdot (2x + 1) =
    2 \cdot (x^2 + 8x + 15) - (2x^2 + x - 14x - 7) =
    2x^2 + 16x + 30 - (2x^2 - 13x - 7) = 2x^2 + 16x + 30 - 2x^2 + 13x + 7 = \mathbf{29x + 37}\) % d)

  \item \({(2x + y)}^2 - {(x - 3y)}^2 = 4x^2 + 4xy + y^2 - (x^2 - 6xy + 9y^2) =
    4x^2 + 4xy + y^2 - x^2 + 6xy - 9y^2 =
    \mathbf{3x^2 + 10xy - 8y^2}\) % e)

  \item \((2a - 3b) \cdot (3a + b) - 2 \cdot (3a + b) \cdot (a - b) =
    6a^2 + 2ab - 9ab - 3b^2 - 2 \cdot (3a^2 - 3ab + ab - b^2) =
    6a^2 - 7ab - 3b^2 - 2 \cdot (3a^2 - 2ab - b^2) = 6a^2 - 7ab - 3b^2 - 6a^2 + 4ab + 2b^2 =
    \mathbf{-3ab - b^2}\) % f)

\end{enumerate}

\subsection{Aufgaben zu Kapitel 3.5 und 3.6 \\
  - Die binomischen Formeln \\
  - Anwendungsbeispiele für die 1. und 2. binomische Formel}

\subsubsection{Benutze die \textit{binomische Formeln}:} % I.

\begin{enumerate}[label=\alph*)]

  \item \({(x + 5)}^2 = \mathbf{x^2 + 10x + 25}\) % a)

  \item \({(x + 2)}^2 = \mathbf{x^2 + 4x + 4}\) % b)

  \item \((x + 3) \cdot (x + 3) = \mathbf{x^2 + 6x + 9}\) % c)

  \item \({(x + y)}^2 = \mathbf{x^2 + 2xy + y^2}\) % d)

  \item \({(x + 3y)}^2 = \mathbf{x^2 + 6xy + 9y^2}\) % e)

  \item \({(5x + 7y)}^2 = \mathbf{25x^2 + 70xy + 49y^2}\) % f)

  \item \({(3u + 2v)}^2 = \mathbf{9u^2 + 12uv + 4v^2}\) % g)

  \item \({(a + 5b)}^2 = \mathbf{a^2 + 10ab + 25b^2}\) % h)

  \item \({(x + 5y)}^2 = \mathbf{x^2 + 10xy + 25y^2}\) % i)

  \item \({(x - 5y)}^2 = \mathbf{x^2 - 10xy + 25y^2}\) % j)

  \item \({(2x - y)}^2 = \mathbf{4x^2 - 4xy + y^2}\) % k)

  \item \({(-2x + 3y)}^2 = \mathbf{4x^2 - 12xy + 9y^2}\) % l)

  \item \({(-x - 7)}^2 = \mathbf{x^2 + 14x + 49}\) % m)

  \item \({(2a - 3b)}^2 = \mathbf{4a^2 - 12ab + 9b^2}\) % n)

  \item \({(7x - 3)}^2 = \mathbf{49x^2 - 42x + 9}\) % o)

\end{enumerate}

\subsubsection{Berechne:} % II.

\begin{enumerate}[label=\alph*)]

  \item \({\big( 5 \cdot (x + 1) \big)}^2 = 25 \cdot (x^2 + 2x + 1) = \mathbf{25x^2 + 50x + 25}\) % a)

  \item \({\big( 4 \cdot (x - 3) \big)}^2 = 16 \cdot (x^2 - 6x + 9) = \mathbf{16x^2 - 96x + 144}\) % b)

  \item \({\big( 3 \cdot (5x + 2) \big)}^2 = 9 \cdot (25x^2 + 20x + 4) = \mathbf{225x^2 + 180x + 36}\) % c)

  \item \({\big( 2y \cdot (3x + 1) \big)}^2 = 4y^2 \cdot (9x^2 + 6x + 1) = \mathbf{36x^2y^2 + 24xy^2 + 4y^2}\) % d)

\end{enumerate}

\subsubsection{Vereinfache mithilfe der \textit{binomischen Formeln}:} % III.

\begin{enumerate}[label=\alph*)]

  \item \({(x + 3)}^2 + {(x - 3)}^2 = x^2 + 6x + 9 + x^2 - 6x + 9 = \mathbf{2x^2 + 18}\) % a)

  \item \({(x + 5)}^2 + {(x - 5)}^2 = x^2 + 10x + 25 + x^2 - 10x + 25 = \mathbf{2x^2 + 50}\) % b)

  \item \({(2x + 1)}^2 - {(3x - 2)}^2 + 3 \cdot {(x + 5)}^2 =
    4x^2 + 4x + 1 - (9x^2 - 12x + 4) + 3 \cdot (x^2 + 10x + 25) =
    4x^2 + 4x + 1 - 9x^2 + 12x - 4 + 3x^2 + 30x + 75 = \mathbf{-2x^2 + 46x + 72}\) % c)

  \item \({(x + 7)}^2 + {(x - 7)}^2 + 2 \cdot (x - 7) \cdot (x + 7) =
    x^2 + 14x + 49 + x^2 - 14x + 49 + 2 \cdot (x^2 - 49) = 2x^2 + 98 + 2x^2 - 98 = \mathbf{4x^2}\) % d)

  \item \({(x^2 + 3)}^2 - {(x^2 - 2)}^2 + 4 \cdot {(x + 3)}^2 =
    x^4 + 6x^2 + 9 - x^4 + 4x^2 - 4 + 4 \cdot (x^2 + 6x + 9) =
    10x^2 + 5 + 4x^2 + 24x + 36 = \mathbf{14x^2 + 24x + 41}\) % e)

\end{enumerate}

\subsubsection{Berechne mithilfe der \textit{binomische Formeln}:} % IV.

\begin{enumerate}[label=\alph*)]

  \item \({101}^2 = {(100 + 1)}^2 = 100^2 + 2 \cdot 100 \cdot 1 + 1^2 = 10.000 + 200 + 1 =
    \mathbf{10.201}\) % a)

  \item \({103}^2 = {(100 + 3)}^2 = 100^2 + 2 \cdot 100 \cdot 3 + 3^2 = 10.000 + 600 + 9 =
    \mathbf{10.609}\) % b)

  \item \({1.008}^2 = {(1.000 + 8)}^2 = 1.000^2 + 2 \cdot 1.000 \cdot 8 + 8^2 =
    1.000.000 + 16.000 + 64 = \mathbf{1.016.064}\) % c)

  \item \({205}^2 = {(200 + 5)}^2 = 200^2 + 2 \cdot 200 \cdot 5 + 5^2 = 40.000 + 2.000 + 25 =
    \mathbf{42.025}\) % d)

  \item \({99}^2 = {(100 - 1)}^2 = 100^2 - 2 \cdot 100 \cdot 1 + 1^2 = 10.000 - 200 + 1 =
    \mathbf{9.801}\) % e)

  \item \({94}^2 = {(100 - 6)}^2 = 100^2 - 2 \cdot 100 \cdot 6 + 6^2 = 10.000 - 1.200 + 36 =
    \mathbf{8.836}\) % f)

  \item \({993}^2 = {(1000 - 7)}^2 = 1000^2 - 2 \cdot 1000 \cdot 7 + 7^2 = 1.000.000 - 14.000 + 49 =
    \mathbf{986.049}\) % g)

  \item \({195}^2 = {(200 - 5)}^2 = 200^2 - 2 \cdot 200 \cdot 5 + 5^2 = 40.000 - 2.000 + 25 =
    \mathbf{38.025}\) % h)

\end{enumerate}

\subsection{Aufgaben zu Kapitel 3.7 \\
  - Anwendungsbeispiele für 3. binomische Formel}

\subsubsection{Benutze die \textit{dritte binomische Formel}:} % I.

\begin{enumerate}[label=\alph*)]

  \item \((x + 5) \cdot (x - 5) = \mathbf{x^2 - 25}\) % a)

  \item \((x + 2) \cdot (x - 2) = \mathbf{x^2 - 4}\) % b)

  \item \((x + 3) \cdot (x - 3) = \mathbf{x^2 - 9}\) % c)

  \item \((x + y) \cdot (x - y) = \mathbf{x^2 - y^2}\) % d)

  \item \((x + 3y) \cdot (x - 3y) = \mathbf{x^2 - 9y^2}\) % e)

  \item \((5x + 7y) \cdot (5x - 7y) = \mathbf{25x^2 - 49y^2}\) % f)

  \item \((3u + 2v) \cdot (3u - 2v) = \mathbf{9u^2 - 4v^2}\) % g)

  \item \((a + 5b) \cdot (a - 5b) = \mathbf{a^2 - 25b^2}\) % h)

  \item \((x + 5y) \cdot (x - 5y) = \mathbf{x^2 - 25y^2}\) % i)

  \item \((2x - 3y) \cdot (2x + 3y) = \mathbf{4x^2 - 9y^2}\) % j)

  \item \((x + 2y) \cdot (x - 2y) = \mathbf{x^2 - 4y^2}\) % k)

  \item \((-2x + 3y) \cdot (-2x - 3y) = \mathbf{4x^2 - 9y^2}\) % l)

  \item \((-x + 7) \cdot (x + 7) = (7 - x) \cdot (7 + x) = \mathbf{49 - x^2}\) % m)

  \item \((2a + 3b) \cdot (-2a + 3b) = (3b + 2a) \cdot (3b - 2a) = \mathbf{9b^2 - 4a^2}\) % n)

  \item \((-7x + 3) \cdot (-7x - 3) = \mathbf{49x^2 - 9}\) % o)

\end{enumerate}

\subsubsection{Berechne:} % II.

\begin{enumerate}[label=\alph*)]

  \item \(5 \cdot (x + 1) \cdot (x - 1) = 5 \cdot (x^2 - 1) = \mathbf{5x^2 - 5}\) % a)

  \item \(4 \cdot (x - 3) \cdot (x + 3) = 4 \cdot (x^2 - 9) = \mathbf{4x^2 - 36}\) % b)

  \item \(3 \cdot (-5x + 2) \cdot (5x + 2) = 3 \cdot (4 - 25x^2) = \mathbf{12 - 75x^2}\) % c)

  \item \(2y \cdot (3x + 1) \cdot (3x - 1) = 2y \cdot (9x^2 - 1) = \mathbf{18x^2y - 2y}\) % d)

\end{enumerate}

\subsubsection{Vereinfache mithilfe der \textit{binomischen Formeln}:} % III.

\begin{enumerate}[label=\alph*)]

  \item \((x + 3) \cdot (x - 3) + (x - 2) \cdot (x + 2) - (x + 1) \cdot (x - 1) =
    x^2 - 9 + x^2 - 4 - x^2 + 1 = \mathbf{x^2 - 12}\) % a)

  \item \({(x + 2)}^2 + {(x - 2)}^2 - (x + 3) \cdot (x - 3) =
    x^2 + 4x + 4 + x^2 - 4x + 4 - x^2 + 9 = \mathbf{x^2 + 17}\) % b)

  \item \({(x + 3)}^2 + {(x - 3)}^2 + (x + 3) \cdot (x - 3) =
    x^2 + 6x + 9 + x^2 - 6x + 9 + x^2 - 9 = \mathbf{3x^2 + 9}\) % c)

  \item \(\cfrac{a^2 - b^2}{a + b} = \cfrac{(a + b) \cdot (a - b)}{(a + b)} = \mathbf{a - b}\) % d)

  \item \(\cfrac{a^2 - b^2}{a - b} = \cfrac{(a + b) \cdot (a - b)}{(a - b)} = \mathbf{a + b}\) % e)

  \item \(\cfrac{{(x - y)}^2}{x^2 - y^2} = \cfrac{(x - y) \cdot (x - y)}{(x + y) \cdot (x - y)} =
    \mathbf{\cfrac{x - y}{x + y}}\) % f)

\end{enumerate}

\subsubsection{Berechne mithilfe der \textit{dritten binomischen Formel}:} % IV.

\begin{enumerate}[label=\alph*)]

  \item \(107 \cdot 93 = (100 + 7) \cdot (100 - 7) = 100^2 - 7^2 = 10.000 - 49 = \mathbf{9.951}\) % a)

  \item \(1.006 \cdot 994 = (1.000 + 6) \cdot (1.000 - 6) = 1.000^2 - 6^2 = 1.000.000 - 36 = \mathbf{999.964}\) % b)

  \item \(101 \cdot 99 = (100 + 1) \cdot (100 - 1) = 100^2 - 1^2 = 10.000 - 1 = \mathbf{9.999}\) % c)

  \item \(28 \cdot 32 = (30 - 2) \cdot (30 + 2) = 30^2 - 2^2 = 900 - 4 = \mathbf{896}\) % d)

\end{enumerate}

\end{document}
